\section{Программа}

\subsection{База знаний}
\begin{listing}[Исходный код базы знаний]{lst:t1-prog}{Prolog}
factorial(N, Result) :- 
    N >= 0,
    fact_tail(N, 1, Result).

fact_tail(0, Acc, Acc) :- !.
fact_tail(N, Acc, Result) :-
    N > 0,
    NewAcc is Acc * N,
    N1 is N - 1,
    fact_tail(N1, NewAcc, Result).


fibonacci(N, Result) :-
    N >= 0,
    fib_tail(N, 0, 1, Result).

fib_tail(0, A, _, A) :- !.
fib_tail(N, A, B, Result) :-
    N > 0,
    NextB is A + B,
    N1 is N - 1,
    fib_tail(N1, B, NextB, Result).
\end{listing}

\subsection{Вопросы задания}

Для получения факториала заданного числа необходимо задать вопрос: 
существует ли такой Результат при конкретизированном Числе, для которого 
удовлетворяется отношение factorial.

\begin{listing}[Поиск факториала числа N]{lst:q1}{Prolog}
?- factorial(N, Res).
\end{listing}

Для получения заданного числа Фибоначчи необходимо задать вопрос: существует
ли такой Результат при конкретизированном Числе, для которого 
удовлетворяется отношение fibonacci.

\begin{listing}[Поиск всех дедушек]{lst:q2}{Prolog}
?- fibonacci(N, Res).
\end{listing}
